\title{DeepSeekMath: Pushing the Limits of Mathematical Reasoning in Open Language Models}

\begin{abstract}
The task of mathematical reasoning remains particularly demanding for language models due to its inherently intricate and systematic nature. In this work, we present DeepSeekMath~7B, an extension of DeepSeek-Coder-Base-v1.5 7B, further pre-trained using 120B math-focused tokens derived from Common Crawl, in conjunction with both natural language and code datasets. DeepSeekMath~7B attains an outstanding 51.7\% on the MATH benchmark at the competition level, all without resorting to supplementary toolkits or majority voting methods, and reaches a performance comparable to that of Gemini-Ultra and GPT-4. With self-consistency computed over 64 generated samples, DeepSeekMath~7B achieves a score of 60.9\% on MATH. The enhanced mathematical reasoning skills of DeepSeekMath~can be credited to two primary components: First, the extensive value present in open-access web resources is exploited via a rigorously developed data selection process. Second, we propose Group Relative Policy Optimization (GRPO), an adapted form of Proximal Policy Optimization (PPO), designed to both advance mathematical reasoning and improve memory efficiency within PPO.
\end{abstract}

\section{Introduction}

Large language models (LLM) have dramatically transformed methodologies for mathematical reasoning within artificial intelligence, leading to notable progress on quantitative reasoning benchmarks \citep{MATH} as well as on geometry reasoning benchmarks \citep{trinh2024solving}. Additionally, these systems have shown substantial utility in supporting humans tackling challenging mathematical tasks \citep{tao}. Nonetheless, state-of-the-art systems like GPT-4 \citep{gpt4} and Gemini-Ultra \citep{gemini} remain closed to the public, while available open-source alternatives still exhibit a substantial performance gap.

This work presents DeepSeekMath, a specialized language model tailored for mathematical reasoning that substantially surpasses existing open-source alternatives and nearly attains GPT-4's proficiency on scholarly evaluations. Our approach hinges on assembling the DeepSeekMath Corpus, a meticulously curated pre-training dataset containing 120B math-specific tokens. To gather this corpus, we rely on Common Crawl (CC) data and employ a classifier based on fastText \citep{joulin2016fasttext}. During the first phase, the classifier is trained using positive samples drawn from OpenWebMath \citep{openwebmath}, with negative samples taken from a varied selection of non-mathematical web content. Leveraging this trained classifier, we identify further positive samples within the CC, subsequently refining these through human annotation. The augmented and annotated data then serve to retrain and enhance the classifier’s effectiveness. According to our assessment, the resultant large-scale corpus maintains exceptional quality: DeepSeekMath-Base 7B achieves a score of 64.2\% on the GSM8K benchmark \citep{gsm8k} and 36.2\% on the MATH competition dataset \citep{MATH}, exceeding the performance of Minerva 540B \citep{minerva}. Moreover, because the DeepSeekMath Corpus encompasses multiple languages, we observe measurable gains in Chinese math benchmarks \citep{wei2023cmath,agieval}. We posit that our practices in mathematical corpus curation may provide a foundation for future research, leaving considerable opportunity for subsequent advancements.

DeepSeekMath-Base builds upon the DeepSeek-Coder-Base-v1.5 7B model \citep{deepseek-coder}, as initializing with a code-oriented pretrained model proves to be more effective than utilizing a general-purpose language model. Additionally, our results show that training on mathematical tasks boosts performance not only in mathematical domains but also yields gains on benchmarks such as MMLU \citep{mmlu} and BBH \citep{bbh}. This suggests that math training enhances both mathematical proficiency and overall reasoning ability.

Following pre-training, we conduct mathematical instruction tuning on DeepSeekMath-Base utilizing chain-of-thought \citep{cot}, program-of-thought \citep{pot,pal}, and tool-integrated reasoning \citep{tora} datasets. The produced model, DeepSeekMath-Instruct 7B, surpasses all other 7B models and achieves performance on par with leading open-source instruction-tuned models of 70B scale.

Additionally, we propose Group Relative Policy Optimization (GRPO), a reinforcement learning (RL) method that builds upon Proximal Policy Optimization (PPO) \citep{schulman2017proximal}. Unlike standard PPO, GRPO eliminates the need for a critic model and instead infers the baseline using group scores, leading to notable savings in computational resources during training. Leveraging only a limited portion of English instruction tuning data, GRPO demonstrates notable gains over the strong DeepSeekMath-Instruct baseline, with marked improvements across both in-domain datasets (GSM8K: 82.9\% $\rightarrow$ 88.2\%, MATH: 46.8\% $\rightarrow$ 51.7\%) and out-of-domain benchmarks such as CMATH (84.6\% $\rightarrow$ 88.8\%) in the reinforcement learning stage. We further present a comprehensive framework for interpreting a range of techniques—including Rejection Sampling Fine-Tuning (RFT) \citep{yuan2023scaling}, Direct Preference Optimization (DPO) \citep{dpo}, PPO, and GRPO—showing that each may be seen as a variant or simplification of RL. Extensive empirical analyses—such as comparisons between online and offline training, outcome-based versus process-based supervision, as well as single-turn versus multi-step RL—are conducted to dissect the core components underlying this unified approach. Finally, we discuss the mechanisms by which RL enhances instruction-tuned model performance and outline promising avenues for future research aimed at advancing RL efficacy within this cohesive framework.

\subsection{Contributions}
We present scalable approaches to mathematical pre-training, in addition to investigating and analyzing the application of reinforcement learning.

\textbf{Math Pre-Training at Scale}

\begin{itemize}[topsep=0pt]
    \item This study demonstrates that Common Crawl, an openly available data source, holds substantial mathematically relevant content. Utilizing a carefully crafted pipeline for selecting data, we curated the DeepSeekMath Corpus, a large-scale dataset comprising 120B tokens from mathematically filtered web pages. This corpus is nearly seven times larger than the math web dataset employed by Minerva \citep{minerva}, and approximately nine times greater than the newly introduced OpenWebMath \citep{openwebmath}.

\item DeepSeekMath-Base 7B, our pre-trained base model, matches the performance level of Minerva 540B \citep{minerva}, suggesting that sheer parameter count is not the sole determinant of proficiency in mathematical reasoning. Notably, even models of a more modest size can perform strongly when trained on carefully curated, high-quality datasets.

\item We present the results of our experiments involving math training. Our data indicate that pretraining models on code before introducing math tasks enhances their performance on mathematical problems, regardless of whether tools are employed. This contributes evidence to the ongoing debate around the question: \textit{does code training enhance reasoning skills?} Our findings support the idea that it does, particularly in the context of mathematical reasoning.

\item While it is a widespread practice to utilize arXiv papers for training, particularly within mathematics-focused research, this strategy does not yield significant gains on any of the mathematical benchmarks considered in this study.
\end{itemize}

\textbf{Exploration and Analysis of Reinforcement Learning}

\begin{itemize}[topsep=0pt]
    \item We present Group Relative Policy Optimization (GRPO), a novel reinforcement learning approach characterized by both efficiency and efficacy. Unlike standard algorithms such as Proximal Policy Optimization (PPO), GRPO omits the use of a critic network and instead derives the baseline directly from group-level scores, which substantially lowers the computational requirements during training.
    \item We show that GRPO yields marked improvements to our instruction-tuned model DeepSeekMath-Instruct when leveraging only instruction-tuning datasets. Notably, we also report improved generalization to out-of-domain tasks throughout the reinforcement learning procedure.
    \item We establish a comprehensive framework for interpreting a range of techniques, including RFT, DPO, PPO, and GRPO. Moreover, we carry out thorough empirical studies—comparing, for instance, online versus offline training, outcome-based versus process-based supervision, and single-step versus multi-iteration reinforcement learning—to deeply examine the key components of this unified framework.
    \item Drawing on insights from our unified paradigm, we analyze the core factors contributing to the success of reinforcement learning methods, and highlight several promising avenues for advancing the effectiveness of reinforcement learning in large language models.
\end{itemize}

\subsection{Summary of Evaluations and Metrics}

\begin{itemize}[topsep=0pt]
    \item \textbf{English and Chinese Mathematical Reasoning}: Our study involves an extensive evaluation of models on both English and Chinese mathematics benchmarks, ranging from elementary to collegiate problem sets. The English suite features GSM8K \citep{gsm8k}, MATH \citep{MATH}, SAT \citep{llemma}, OCW Courses \citep{minerva}, and MMLU-STEM \citep{mmlu}. The Chinese suite encompasses MGSM-zh \citep{mgsm}, CMATH \citep{wei2023cmath}, Gaokao-MathCloze \citep{agieval}, and Gaokao-MathQA \citep{agieval}. Our evaluation methodology includes testing whether models can produce complete written solutions unaided by external tools, as well as assessing their proficiency in solving tasks with the use of Python.

On English evaluation benchmarks, DeepSeekMath-Base achieves results comparable to the proprietary Minerva 540B model \citep{minerva}, and outperforms all existing open-source base models such as Mistral 7B \citep{mistral} and Llemma-34B \citep{llemma}, irrespective of their use of math-oriented pre-training—often by a notable margin. Particularly, DeepSeekMath-Base exhibits superior performance on Chinese benchmarks, which may be attributed to our divergence from previous methodologies \citep{minerva,llemma} that rely solely on English math pre-training datasets, as our approach incorporates high-quality data in other languages. Following instruction tuning specific to mathematics and the application of reinforcement learning, the advanced models DeepSeekMath-Instruct and DeepSeekMath-RL attain robust results, becoming the first open-source models to achieve over 50\% accuracy on the competition-grade MATH dataset.

\item \textbf{Formal Mathematics}: To assess DeepSeekMath-Base, we employ the informal-to-formal theorem proving benchmark introduced by \citep{dsp_proof}, utilizing miniF2F \citep{minif2f} with the Isabelle proof assistant \citep{isabelle}. Results indicate that DeepSeekMath-Base achieves notable few-shot performance in automatic formalization tasks.
\item \textbf{Natural Language Understanding, Reasoning, and Code}: For a holistic evaluation of the model’s abilities in general comprehension, logical reasoning, and programming, DeepSeekMath-Base is benchmarked on the Massive Multitask Language Understanding (MMLU) suite \citep{mmlu}, which spans 57 diverse multiple-choice domains; BIG-Bench Hard (BBH) \citep{bbh}, containing 23 demanding tasks requiring complex, multi-step inference; and two standard code evaluation datasets, HumanEval \citep{codex} and MBPP \citep{mbpp}. The results suggest that pre-training on mathematical data enhances the model’s performance in both linguistic and reasoning benchmarks.

\section{Math Pre-Training}

\subsection{Data Collection and Decontamination} This section details the methodology used to assemble the DeepSeekMath Corpus utilizing Common Crawl data. Figure \ref{fig:data_collect} illustrates an iterative workflow, which systematically extracts a substantial mathematical corpus from Common Crawl, beginning with a seed collection comprised of a limited yet high-quality set of mathematics-focused data. Importantly, this methodology is flexible and can be extended to other areas, including code.

Our process begins by selecting OpenWebMath \citep{openwebmath}, a curated set of high-quality mathematical materials from the web, to serve as our seed dataset. Based on this corpus, we develop a fastText model \citep{joulin2016fasttext} aimed at retrieving additional mathematical web pages similar in nature to those in OpenWebMath. For the construction of our training data, we sample 500,000 entries from the seed corpus as positive examples and an equivalent number of web pages from Common Crawl as negative examples. Model training is carried out with a publicly available fastText implementation\footnote{\url{https://fasttext.cc}}, where hyperparameters are set as follows: vector size at 256, learning rate of 0.1, a maximum n-gram length of 3, minimum word frequency threshold of 3, and the number of epochs set to 3. To curtail the extent of Common Crawl, we conduct both URL-based deduplication and near-duplicate removal, resulting in a reduced dataset of approximately 40 billion HTML documents. Using the trained fastText model, we retrieve mathematical web content from this deduplicated collection. In order to eliminate low-quality material, the recalled web pages are ordered based on their fastText model scores, with only the highest-scoring content being retained. The specific size of the preserved dataset is further determined through pre-training experiments with token cutoffs at 40B, 80B, 120B, and 160B. For the initial stage, we opt to retain the top 40B tokens.

Following the initial round of data gathering, a significant number of mathematical web pages remain unretrieved. This shortfall is chiefly attributed to the limited diversity in the positive samples used to train the fastText model. To address this, we expand our list of mathematical web sources to further enhance the seed dataset, facilitating better model optimization. Our strategy begins with segmenting the entire Common Crawl into non-overlapping domains, where each domain consists of all web pages under a common base URL. We then compute the proportion of web pages from each domain that were acquired during the first collection phase. If more than 10\% of a domain's web pages are included in this set, we categorize the domain as math-related (for instance, \url{mathoverflow.net}).

Next, we manually tag URLs that lead to mathematical content within these selected domains (such as \url{mathoverflow.net/questions}). Web pages linked to these URLs that were not previously collected are incorporated into the seed corpus. Through this methodology, we broaden the pool of positive examples, enabling us to train an improved fastText model with a higher capability to recall mathematical content in the next iteration. After performing four rounds of data collection, we accumulate 35.5M mathematical web pages, corresponding to 120B tokens. By the fourth iteration, analysis reveals that roughly 98\% of the corpus had already been gathered by the previous round, prompting us to halt further data collection.

To mitigate the risk of benchmark contamination, we adopt the methodology outlined in \cite{deepseek-coder}, filtering out any web pages that include questions or answers from prominent English mathematical benchmarks such as GSM8K \citep{gsm8k} and MATH \citep{MATH}, as well as from Chinese resources like CMATH \citep{wei2023cmath} and AGIEval \citep{agieval}. Our approach involves removing any segment from the mathematical training dataset if it contains an exact match of a 10-gram sequence found in these evaluation benchmarks. For benchmark passages shorter than 10 grams but at least 3 grams in length, exact substring matching is similarly utilized to eliminate potentially contaminated sources.

\subsection{Validating the Quality of the DeepSeekMath~Corpus}

We run pre-training experiments to investigate how the DeepSeekMath~Corpus is compared with the recently released math-training corpora:

\begin{itemize}[topsep=0pt]
    \item \textbf{MathPile} \citep{mathpile}: This dataset comprises 8.9B tokens and is formed by collecting data from various sources such as textbooks, Wikipedia, ProofWiki, CommonCrawl, StackExchange, and arXiv, with arXiv providing more than 85\% of the content;
    \item \textbf{OpenWebMath} \citep{openwebmath}: Drawn from CommonCrawl, this corpus has been curated to include only materials relevant to mathematics, reaching a size of 13.6B tokens;
    \item \textbf{Proof-Pile-2} \citep{llemma}: Composed of mathematical documents from OpenWebMath, AlgebraicStack (providing 10.3B tokens of mathematics-related code), and arXiv papers (28.0B tokens). Consistent with the approach in \cite{llemma}, we employ a data mixture ratio of 2:4:1 for arXiv:Web:Code during experiments with Proof-Pile-2.
\end{itemize}

\subsubsection{Training Setting}
\label{sec:quality-policy}
We conduct mathematical instruction tuning on a broadly pre-trained language model comprising 1.3B parameters, following the DeepSeek LLM framework \citep{deepseek-llm}, referred to here as DeepSeek-LLM 1.3B. For each individual mathematical dataset, the model is trained separately, processing 150B tokens in total. All experimental runs utilize the lightweight HAI-LLM \citep{haillm} training infrastructure to ensure efficiency. Adhering to the DeepSeek LLM training protocols, we adopt the AdamW optimizer \citep{adamW} with settings $\beta_1=0.9$, $\beta_2=0.95$, and $\mathrm{weight\_decay}=0.1$. The learning rate follows a multi-phase schedule: it ramps up to its maximum after 2,000 steps of warmup, then reduces to 31.6\% of this peak value at 80\% of total training steps, and further declines to 10.0\% by the 90\% mark of training completion. The peak learning rate is fixed at 5.3e-4, and training utilizes a batch size accommodating 4M tokens and a sequence context length of 4K.

\subsubsection{Evaluation Results}

\textbf{The DeepSeekMath~Corpus is of high quality, covers multilingual mathematical content, and is the largest in size.}

\begin{itemize}[topsep=0pt]
    \item \textbf{High-quality}:
    We assess downstream task performance across eight mathematical benchmarks employing few-shot chain-of-thought prompting \cite{cot}. According to Table \ref{tab:corpora_comparison}, models trained on the DeepSeekMath~Corpus consistently surpass others in accuracy. As illustrated in Figure \ref{fig:corpora_comparisons}, a model utilizing the DeepSeekMath Corpus outperforms the Proof-Pile-2 benchmark at 50B tokens (equivalent to one complete epoch of Proof-Pile-2), suggesting superior average data quality within the DeepSeekMath Corpus. 
    \item \textbf{Multilingual}:
    The DeepSeekMath~Corpus incorporates a variety of languages, with English and Chinese being the most prevalent. Table \ref{tab:corpora_comparison} demonstrates that models trained with the DeepSeekMath~Corpus attain higher mathematical reasoning scores for both English and Chinese tasks. In comparison, prevailing mathematical corpora, which are predominantly English-focused, provide minimal gains and sometimes even impair Chinese mathematical reasoning abilities.
    \item \textbf{Large-scale}:
    The DeepSeekMath~Corpus boasts a scale multiple times greater than that of previous mathematical corpora. Figure \ref{fig:corpora_comparisons} reveals that DeepSeek-LLM 1.3B, after being trained on the DeepSeekMath~Corpus, achieves a sharper learning trajectory and more persistent enhancements. In stark contrast, baseline corpora are considerably smaller, have undergone several epochs during training, and their associated model performance levels off much sooner.
\end{itemize}

\subsection{Training and Evaluating DeepSeekMath-Base 7B}

This section details the DeepSeekMath-Base 7B, a foundational model notable for its advanced reasoning capabilities in mathematics. The model is built upon DeepSeek-Coder-Base-v1.5 7B \citep{deepseek-coder} and undergoes training over 500B tokens. The composition of the training data includes 56\% from the DeepSeekMath Corpus, 4\% from AlgebraicStack, 10\% sourced from arXiv, 20\% consisting of Github code, with the remaining 10\% comprising natural language content in both English and Chinese drawn from Common Crawl. We generally employ the training procedures outlined in Section \ref{sec:quality-policy}, with the exception that the peak learning rate is adjusted to 4.2e-4 and the batch size is configured to 10M tokens.

Our evaluation thoroughly investigates the mathematical proficiency of DeepSeekMath-Base 7B, emphasizing its capacity to independently generate complete mathematical solutions, utilize tools to tackle mathematical problems, and engage in formal theorem proving. In addition to mathematics, we expand our analysis to offer a broader characterization of the base model, examining its abilities in natural language understanding, reasoning, and programming.

\paragraph{Mathematical Problem Solving with Step-by-Step Reasoning}

We assess the ability of DeepSeekMath-Base to solve mathematical problems by employing few-shot chain-of-thought prompting \citep{cot} across eight distinct benchmarks in both English and Chinese. These evaluation benchmarks span quantitative reasoning tasks, such as GSM8K \citep{gsm8k}, MATH \citep{MATH}, and CMATH \citep{wei2023cmath}, as well as multiple-choice assessments including MMLU-STEM \citep{mmlu} and Gaokao-MathQA \citep{agieval}. The selection encompasses a wide spectrum of mathematical domains, ranging from elementary to university-level topics.

As indicated in Table \ref{tab:base_math_cot}, DeepSeekMath-Base 7B demonstrates superior results across all eight evaluated benchmarks when compared with other open-source base models, such as the prevalent Mistral 7B \citep{mistral} and the more recent Llemma 34B \citep{llemma}, the latter of which benefits from targeted math training using Proof-Pile-2 \citep{llemma}. Particularly on the challenging MATH dataset, which features competition-level problems, DeepSeekMath-Base exceeds all currently available open-source base models by more than 10\% absolute and even outperforms the closed-source Minerva 540B \citep{minerva}—a model that is 77 times larger, based on PaLM \citep{palm}, and subjected to extensive mathematical text training.

\paragraph{Mathematical Problem Solving with Tool Use} We assess the capability of models in program-assisted mathematical reasoning on the GSM8K and MATH benchmarks by employing few-shot program-of-thought prompting methods \citep{pot,pal}. For each task, models generate a Python script to find a solution, making use of libraries such as \textit{math} and \textit{sympy} to handle complex calculations. The output produced by running the code is taken as the model’s answer. As summarized in Table \ref{tab:base_math_pot_proof}, DeepSeekMath-Base 7B surpasses the performance of the previous state-of-the-art, Llemma 34B.

\paragraph{Formal Mathematics}

The automation of formal proofs offers significant advantages in promoting the precision and dependability of mathematical arguments while also improving productivity, and has garnered heightened interest recently. In this study, we assess the performance of DeepSeekMath-Base 7B on the informal-to-formal proof task as introduced in \citep{dsp_proof}. This task involves producing a formal proof given an informal description, its corresponding formal representation, and an informal proof. Our evaluation utilizes the miniF2F benchmark \citep{minif2f}, which focuses on formal proofs for Olympiad-level mathematics. For each task, we prompt the model with a few examples and require it to generate a formal proof in Isabelle. Consistent with the methodology of \cite{dsp_proof}, we task the models with composing proof sketches, and then employ the automated prover Sledgehammer \citep{sledgehammer} to complete the detailed reasoning steps. The results, summarized in Table \ref{tab:base_math_pot_proof}, indicate that DeepSeekMath-Base 7B performs robustly in the context of proof autoformalization.

\paragraph{Natural Language Understanding, Reasoning, and Code}

We assess the natural language understanding abilities of the model using MMLU \citep{mmlu}, its reasoning skills with BBH \citep{bbh}, and its coding performance with HumanEval \citep{codex} and MBPP \citep{mbpp}. Table \ref{tab:understanding_reasoning_code} demonstrates that DeepSeekMath-Base 7B achieves considerable improvements over its predecessor, DeepSeek-Coder-Base-v1.5 \citep{deepseek-coder}, on both MMLU and BBH, highlighting the beneficial effects of mathematical pretraining on both language understanding and reasoning tasks. Furthermore, the inclusion of code tokens during continued pretraining allows DeepSeekMath-Base 7B to largely retain the coding performance of DeepSeek-Coder-Base-v1.5 across both code-focused benchmarks. In summary, DeepSeekMath-Base 7B consistently surpasses the more general Mistral 7B \citep{mistral} when evaluated on all three reasoning and coding tasks.

\section{Supervised Fine-Tuning}

\subsection{SFT Data Curation}

We develop a comprehensive mathematical instruction-tuning dataset encompassing both English and Chinese questions drawn from a range of mathematical disciplines and exhibiting diverse degrees of complexity. Each problem is accompanied by a solution in chain-of-thought (CoT) \citep{cot}, program-of-thought (PoT) \citep{pot,pal}, or tool-integrated reasoning format \citep{tora}. The dataset includes a total of 776K training examples.

\begin{itemize}[topsep=0pt]
    \item \textbf{English mathematical datasets}: We provide tool-assisted solution annotations for problems from GSM8K and MATH, and include a selection of examples from MathInstruct \citep{MathInstruct} as well as the training split of Lila-OOD \citep{lila}, both featuring solutions developed using either CoT or PoT strategies. The English dataset encompasses a broad spectrum of mathematical disciplines, including but not limited to algebra, probability, number theory, calculus, and geometry.
    \item \textbf{Chinese mathematical datasets}: Our collection consists of Chinese K-12 math questions, covering 76 distinct sub-domains like linear equations, each with answers annotated in both CoT and tool-augmented reasoning styles.
\end{itemize}

\subsection{Training and Evaluating DeepSeekMath-Instruct 7B}

In this section, we present DeepSeekMath-Instruct 7B, which is derived from DeepSeekMath-Base through instruction tuning with mathematical tasks. During training, we randomly concatenate samples until the sequence reaches the maximum allowed context window of 4K tokens. The training process consists of 500 steps, utilizing a batch size of 256 and maintaining a fixed learning rate of 5e-5 throughout.

We evaluate models' mathematical performance both without and with tool use, on 4 quantitative reasoning benchmarks in English and Chinese. We benchmark our model against the leading models of the time:

\begin{itemize}[topsep=0pt]
    \item \textbf{Closed-source models} encompass the following: (1) the GPT series, with GPT-4 \citep{gpt4} and GPT-4 Code Interpreter~\footnote{\url{https://openai.com/blog/chatgpt-plugins\#\#code-interpreter}} being the most advanced; (2) Gemini Ultra and Pro \citep{gemini}; (3) Inflection-2 \citep{inflection-2}; (4) Grok-1~\footnote{\url{https://x.ai/model-card}}; in addition, several recent models from Chinese firms such as (5) Baichuan-3~\footnote{\url{https://www.baichuan-ai.com}}; and (6) the most recent GLM-4~\footnote{\url{https://open.bigmodel.cn/dev/api\#glm-4}}.

    from the GLM family \citep{glm}.

These models serve broad applications and typically have been subjected to several alignment strategies.
\item \textbf{Open-source models} consist of: standard large-scale models such as (1) DeepSeek-LLM-Chat 67B \citep{deepseek-llm}, (2) Qwen 72B \citep{qwen}, (3) SeaLLM-v2 7B \citep{seallm}, and (4) ChatGLM3 6B \citep{chatglm3}, along with math-focused variants including (5) InternLM2-Math 20B~\footnote{\url{https://github.com/InternLM/InternLM-Math}}—which enhances InternLM2 via targeted math instruction and subsequent instruction tuning, (6) Math-Shepherd-Mistral 7B, leveraging PPO training \citep{schulman2017proximal} on Mistral 7B \citep{mistral} guided by a process-supervised reward model, (7) the WizardMath series \citep{wizardmath}, advancing mathematical reasoning skills in Mistral 7B and Llama-2 70B \citep{llama2} through evolve-instruct (a type of instruction tuning based on AI-generated instructions) and PPO optimization, primarily training on data sets such as GSM8K and MATH, (8) MetaMath 70B \citep{metamath}, in which Llama-2 70B undergoes fine-tuning on a broadened version of GSM8K and MATH, (9) ToRA 34B \cite{tora}, derived from CodeLlama 34B and optimized for tool-assisted mathematical reasoning, and (10) MAmmoTH 70B \citep{MathInstruct}, produced by instruction-tuning Llama-2 70B on MathInstruct.

Table \ref{tab:sft_rl_math} illustrates that, when tool use is not permitted during evaluation, DeepSeekMath-Instruct 7B exhibits robust step-by-step reasoning abilities. Specifically, the model achieves higher accuracy on the competition-level MATH benchmark than all available open-source models and most closed-source systems (including Inflection-2 and Gemini Pro), exceeding their performance by a margin of at least 9\% absolute. This advantage holds even against models with substantially larger parameter counts (such as Qwen 72B) or those that have undergone dedicated reinforcement learning with a focus on mathematical proficiency (for instance, WizardMath-v1.1 7B). Although DeepSeekMath-Instruct matches the performance of proprietary Chinese models like GLM-4 and Baichuan-3 on the MATH dataset, it continues to lag behind leading models such as GPT-4 and Gemini Ultra.

In an evaluation scenario permitting the combination of natural language reasoning with programmatic tool utilization, DeepSeekMath-Instruct 7B attains nearly 60\% accuracy on the MATH benchmark, outperforming all prior open-source models. Across additional benchmarks, our model performs on par with DeepSeek-LLM-Chat 67B, the previous state-of-the-art model despite it being an order of magnitude larger.  
\section{Reinforcement Learning}

\subsection{Group Relative Policy Optimization} The use of reinforcement learning (RL) has shown strong potential in enhancing the mathematical reasoning skills of LLMs following the Supervised Fine-Tuning (SFT) phase \citep{wang2023math,wizardmath}. Here, we present Group Relative Policy Optimization (GRPO), a novel RL approach designed to be both efficient and effective.

\subsubsection{From PPO to GRPO} Proximal Policy Optimization (PPO) \citep{schulman2017proximal} is an actor-critic RL algorithm that is widely used in the RL fine-tuning stage of LLMs \citep{ouyang2022training}. In particular, it optimizes LLMs by maximizing the following surrogate objective:

\begin{equation}
\footnotesize
    \mathcal{J}_{PPO}(\theta) = \mathbb{E}{[q \sim P(Q),\ o \sim \pi_{\theta_{old}}(O|q)]} \frac{1}{|o|} \sum_{t=1}^{|o|} \min \left[ \frac{\pi_\theta(o_{t} | q, o_{<t})}{\pi_{\theta_{old}}(o_{t} | q, o_{<t})} A_{t},\ \text{clip} \left( \frac{\pi_\theta(o_{t} | q, o_{<t})}{\pi_{\theta_{old}}(o_{t} | q, o_{<t})},\ 1 - \varepsilon,\ 1 + \varepsilon \right) A_{t} \right],
\end{equation}

Here, $\pi_{\theta}$ denotes the current policy while $\pi_{\theta_{old}}$ corresponds to the previous policy version. The variables $q$ and $o$ represent questions and outputs, sampled from the question dataset and generated by the former policy $\pi_{\theta_{old}}$, respectively. The hyper-parameter $\varepsilon$ is employed for clipping within PPO, serving to maintain training stability. The term $A_t$ indicates the advantage, which is estimated using Generalized Advantage Estimation (GAE) \citep{gae}, leveraging rewards $\{r_{\ge t}\}$ alongside predictions from the value function $V_{\psi}$. As a result, PPO necessitates the concurrent training of a value function with the policy network. In order to control reward model over-optimization, the prevalent methodology appends a per-token KL divergence penalty—computed relative to a reference model—to the reward at each time step \citep{ouyang2022training}; that is,

\begin{equation}
    r_{t} = r_\varphi(q, o_{\le t}) - \beta \log\frac{\pi_{\theta}(o_{t}|q, o_{<t})}{\pi_{ref}(o_{t}|q, o_{<t})},
\label{eq:PPO-reward}
\end{equation}

Here, $r_\varphi$ represents the reward model, $\pi_{ref}$ denotes the reference model—commonly the initial SFT model—and $\beta$ stands for the coefficient applied to the KL divergence penalty.

As the value function employed in PPO is typically another model of comparable size as the policy model, it brings a substantial memory and computational burden. Additionally, during RL training, the value function is treated as a baseline in the calculation of the advantage for variance reduction. While in the LLM context, usually only the last token is assigned a reward score by the reward model, which may complicate the training of a value function that is accurate at each token. To address this, as shown in Figure \ref{fig:grpo}, we propose Group Relative Policy Optimization (GRPO), which obviates the need for additional value function approximation as in PPO, and instead uses the average reward of multiple sampled outputs, produced in response to the same question, as the baseline. More specifically, for each question $q$, GRPO samples a group of outputs $\{o_1, o_2, \cdots, o_G\}$  from the old policy  $\pi_{\theta_{old}}$  and then optimizes the policy model by maximizing the following objective:

\begin{equation}
\footnotesize
\begin{split}
    \mathcal{J}_{GRPO}(\theta) &= \mathbb{E}{[q \sim P(Q), \{o_i\}_{i=1}^G \sim \pi_{\theta_{old}}(O|q)]}  \\
    & \frac{1}{G}\sum_{i=1}^G\frac{1}{|o_i|} \sum_{t=1}^{|o_i|} \left\{ \min \left[ \frac{\pi_\theta(o_{i,t} | q, o_{i,<t})}{\pi_{\theta_{old}}(o_{i,t} | q, o_{i,<t})} \hat{A}_{i,t}, \text{clip} \left( \frac{\pi_\theta(o_{i,t} | q, o_{i,<t})}{\pi_{\theta_{old}}(o_{i,t} | q, o_{i,<t})}, 1 - \varepsilon, 1 + \varepsilon \right)  \hat{A}_{i,t} \right] - \beta \mathbb{D}_{KL}\left[\pi_{\theta} || \pi_{ref}\right]\right\} ,
\end{split}
\label{eq:GRPO-obj}
\end{equation}

where $\varepsilon$ and $\beta$ are hyper-parameters, and $\hat{A}_{i,t}$  is the advantage calculated based on relative rewards of the outputs inside each group only, which will be detailed in the following subsections. The group relative way that GRPO leverages to calculate the advantages, aligns well with the comparative nature of rewards models,  as reward models are typically trained on datasets of comparisons between outputs on the same question. Also note that, instead of adding KL penalty in the reward, GRPO regularizes by directly adding the KL divergence between the trained policy and the reference policy to the loss, avoiding complicating the calculation of  $\hat{A}_{i,t}$. And different from the KL penalty term used in (\ref{eq:PPO-reward}), we estimate the KL divergence with the following unbiased estimator \citep{kl_approx}:

\begin{equation}
\small
    \mathbb{D}_{KL}\left[\pi_{\theta} || \pi_{ref}\right] = \frac{\pi_{ref}(o_{i,t}|q,o_{i,<t})}{\pi_{\theta}(o_{i,t}|q,o_{i,<t})} - \log\frac{\pi_{ref}(o_{i,t}|q,o_{i,<t})}{\pi_{\theta}(o_{i,t}|q,o_{i,<t})} - 1,
\end{equation}

which is guaranteed to be positive.

\begin{algorithm}[t]
  \small
  \caption{Iterative Group Relative Policy Optimization}
  \textbf{Input} initial policy model $\pi_{\theta_{\text{init}}}$; reward models $r_\varphi$; task prompts $\mathcal{D}$; 
  hyperparameters $\varepsilon$, $\beta$, $\mu$
  \begin{algorithmic}[1]
    \State Set policy model $\pi_\theta$ to $\pi_{\theta_{\text{init}}}$
    \For{iteration = 1, \dots, I}
       \State Assign reference model $\pi_{ref}$ as $\pi_{\theta}$
      \For{step = 1, \dots, M}
      \State Draw a mini-batch $\mathcal{D}_b$ from $\mathcal{D}$
      \State Assign the previous policy model: $\pi_{\theta_{old}} \gets \pi_{\theta}$ 
      \State For each prompt $q \in \mathcal{D}_b$, generate $G$ outputs $\{o_i\}_{i=1}^G$ sampled from $\pi_{\theta_{old}} (\cdot \mid q)$
      \State Calculate rewards $\{r_i\}_{i=1}^{G}$ for each output $o_i$ by evaluating them with $r_{\varphi}$
      \State Estimate group relative advantages $\hat{A}_{i,t}$ for each token $t$ of $o_i$
      \For{GRPO iteration = 1, \dots, $\mu$}
        \State Adjust the policy model $\pi_{\theta}$ by optimizing the GRPO objective (see Equation \ref{eq:GC-GRPO})
      \EndFor
    \EndFor \State Perform continued training of $r_\varphi$ using a replay strategy. 
    \EndFor 
  \end{algorithmic}
  \textbf{Output} $\pi_\theta$
  \label{alg:iter-grpo}
\end{algorithm}

\subsubsection{Outcome Supervision RL with GRPO} In this framework, given a question $q$, a collection of $G$ outputs $\{o_1, o_2, \cdots, o_G\}$ are drawn from the prior policy $\pi_{\theta_{old}}$. Each output is then evaluated by a reward model, which assigns a set of $G$ corresponding scores $\mathbf{r}=\{r_1, r_2, \cdots, r_G\}$. These rewards are standardized within the group by deducting the mean and dividing by the standard deviation of the set. Under outcome supervision, the normalized reward is assigned as a terminal reward for every output $o_i$, with the computed normalized score used as the advantage $\hat{A}_{i, t}$ for all tokens in that output—specifically, $\hat{A}_{i, t} = \widetilde{r}_i = \frac{r_i- {\rm mean}(\mathbf{r})}{{\rm std}(\mathbf{r})}$. The policy is then updated by maximizing the objective function described in equation (\ref{eq:GRPO-obj}).

\subsubsection{Process Supervision RL with GRPO}  
While outcome supervision restricts feedback to only a final reward after each output, such a strategy may prove inadequate and inefficient for guiding the policy on challenging mathematical reasoning tasks. To address this, and in line with \cite{wang2023math}, we investigate process supervision, where rewards are provided after every intermediate reasoning step. Specifically, given a question $q$ and a set of $G$ sampled output sequences $\{o_1, o_2, \cdots, o_G\}$, a process reward model evaluates each intermediate step, generating a set of stepwise rewards: $\mathbf{R} = \{ \{r_1^{index(1)},\cdots,r_1^{index(K_1)}\}, \cdots,  \{r_G^{index(1)},\cdots,r_G^{index(K_G)}\} \}$. Here, $index(j)$ marks the token position of the $j$-th reasoning step's conclusion, and $K_i$ represents the step count in the $i$-th output. To facilitate stable learning, these rewards are normalized by subtracting their mean and dividing by their standard deviation: $\widetilde{r}_i^{index(j)} = \frac{r_i^{index(j)} - {\rm mean(\mathbf{R})}}{{\rm std(\mathbf{R})}}$.

For each token, the process supervision approach determines its advantage by aggregating the normalized rewards from all subsequent steps: $\hat{A}_{i, t} = \sum_{index(j) \ge t} \widetilde{r}_i^{index(j)}$. The policy is then improved by maximizing the GRPO objective function as specified in equation (\ref{eq:GRPO-obj}).

\subsubsection{Iterative RL with GRPO} During reinforcement learning, the initial reward model may become inadequate for effectively guiding the updated policy model over time. To address this, we implement iterative RL with GRPO. As illustrated in Algorithm \ref{alg:iter-grpo}, this approach involves constructing updated training datasets for the reward model using samples generated by the latest policy model, and then fine-tuning the existing reward model through a replay strategy that integrates 10\% of earlier data. Subsequently, we designate the current policy model as the reference, and proceed to further optimize the policy model with the updated reward model.

\subsection{Training and Evaluating DeepSeekMath-RL}

Our reinforcement learning experiments utilize DeepSeekMath-Instruct 7B as the base model. For RL training, we focus on questions presented in a chain-of-thought style drawn from the GSM8K and MATH subsets of the SFT data, totaling approximately 144,000 problems. We deliberately omit other SFT examples in order to assess how RL affects performance on benchmarks for which data is absent during the RL process. The reward model training set is constructed according to the approach described by \citep{wang2023math}. The initial reward model is trained on DeepSeekMath-Base 7B with a learning rate of $2 \times 10^{-5}$. For the GRPO procedure, we employ a policy learning rate of $1 \times 10^{-6}$. The KL divergence coefficient is assigned a value of 0.04. Each prompt generates $64$ candidate completions. Sequence length and training batch size are both limited to 1024. The policy model undergoes just a single update after each exploration period. We assess DeepSeekMath-RL 7B on the same suite of benchmarks as those used for DeepSeekMath-Instruct 7B. For DeepSeekMath-RL 7B, tasks derived from GSM8K and MATH with chain-of-thought solutions are treated as in-domain, while remaining benchmarks serve as out-of-domain tasks.

Table \ref{tab:sft_rl_math} presents a comparative analysis of open-source and closed-source models featuring both chain-of-thought and tool-augmented reasoning across English and Chinese evaluation datasets. Our key observations include: 1) Employing chain-of-thought reasoning, DeepSeekMath-RL 7B achieves 88.2\% and 51.7\% accuracy on GSM8K and MATH, respectively. These results not only outperform all other open-source models within the 7B–70B parameter range but also exceed the accuracy of most closed-source competitors. 2) Notably, DeepSeekMath-RL 7B has been exclusively trained using chain-of-thought instruction tuning data drawn from GSM8K and MATH, with initialization from DeepSeekMath-Instruct 7B. Although trained on a narrowly defined dataset, it consistently surpasses DeepSeekMath-Instruct 7B across all benchmarks, highlighting the benefits of reinforcement learning in this context.

\section{Discussion}
This section presents an overview of our observations obtained from the pre-training and reinforcement learning (RL) experiments.

\subsection{Lessons Learnt in Pre-Training}

To begin, we present insights gained from the pre-training process. Except where indicated, the training configurations described in Section \ref{sec:quality-policy} are consistently applied. Additionally, throughout this section, references to the DeepSeekMath Corpus pertain to an 89B-token dataset obtained during the second round of data collection.

\subsubsection{Code Training Benefits Mathematical Reasoning}

An often-cited but empirically unproven theory posits that code training enhances reasoning skills. In this work, we present partial evidence related to this claim in the context of mathematics, demonstrating that code training bolsters models' capacity for mathematical reasoning, irrespective of whether tools are employed.

To study how code training affects mathematical reasoning, we experimented with the following two-stage training and one-stage training settings:

\textbf{Two-Stage Training}

\begin{itemize}[topsep=0pt]
    \item \textbf{Code Training for 400B Tokens $\rightarrow$ Math Training for 150B Tokens}: The DeepSeek-LLM 1.3B model is initially trained on 400B code tokens, with subsequent fine-tuning on an additional 150B tokens consisting of mathematical content;
    \item \textbf{General Training for 400B Tokens $\rightarrow$ Math Training for 150B Tokens}: To provide a baseline comparison, we also conduct training where the initial 400B tokens are sampled from a diverse, large-scale general-purpose dataset curated by DeepSeek-AI rather than from code, followed by the same math-centric training. This setup is designed to assess whether code tokens confer specific benefits for mathematical reasoning over general domain data.
\end{itemize}

\textbf{One-Stage Training}

\begin{itemize}[topsep=0pt]
    \item \textbf{Math-Focused Training Using 150B Tokens}: We conduct math-specific training on DeepSeek-LLM 1.3B utilizing a dataset comprised of 150B math tokens;
    \item \textbf{Combined Training with 400B Code Tokens and 150B Math Tokens}: Sequentially applying math training after code training has been found to impair coding abilities. To address this, we explore whether integrating code and math tokens into a unified training phase could both bolster mathematical reasoning skills and help mitigate catastrophic forgetting.
\end{itemize}

\paragraph{Results}
The downstream outcomes across various training configurations are presented in Table \ref{tab:code-math} and Table \ref{tab:code-math-general-eval}.

Incorporating code training improves program-aided mathematical reasoning in both two-stage and one-stage training paradigms. According to Table \ref{tab:code-math}, code training on its own already produces notable gains in solving GSM8K and MATH tasks with Python during the two-stage approach. Introducing math-focused training during the second stage brings additional performance boosts. Notably, in the one-stage training scheme, jointly training on code tokens and math tokens not only addresses the catastrophic forgetting observed in the two-stage method but also fosters enhanced capabilities in both general coding (Table \ref{tab:code-math-general-eval}) and program-aided mathematical reasoning (Table \ref{tab:code-math}).

Engaging in code training leads to enhancements in mathematical reasoning even in the absence of tool assistance. When employing a two-stage training approach, the preliminary phase focused on code yields moderate gains on its own. Furthermore, it facilitates more efficient learning during the following mathematical training phase, culminating in superior results overall. In contrast, merging code tokens and math tokens within a single-stage training framework diminishes mathematical reasoning abilities without tool support. This outcome may be attributed to the relatively small scale of DeepSeek-LLM 1.3B, which likely lacks sufficient capacity to effectively integrate code and mathematical datasets at the same time.

\subsubsection{ArXiv Papers Seem Ineffective in Improving Mathematical Reasoning}

ArXiv papers are commonly included as a component of math pre-training data \citep{minerva,gpt-f,llemma,mathpile}. However, detailed analysis regarding their impact on mathematical reasoning has not been extensively conducted. Perhaps counter-intuitively, according to our experiments, arXiv papers seem ineffective in improving mathematical reasoning. We experiment with models of different sizes, including DeepSeek-LLM 1.3B and DeepSeek-Coder-Base-v1.5 7B \citep{deepseek-coder}, using arXiv corpora that underwent varied processing pipelines:

\begin{itemize}[topsep=0pt]
    \item \textbf{MathPile} \citep{mathpile}: This dataset consists of 8.9 billion tokens, curated through heuristic-based cleaning and filtering, with scientific arXiv articles comprising more than 85\% of its content;
    \item \textbf{ArXiv-RedPajama} \citep{redpajama}: This dataset encompasses all arXiv LaTeX source files after eliminating preambles, macros, comments, and bibliography sections, resulting in a collection of 28.0 billion tokens.
\end{itemize}

For our experimental setup, we independently pretrain DeepSeek-LLM 1.3B on 150B tokens and DeepSeek-Coder-Base-v1.5 7B on 40B tokens, using each arXiv dataset. The results suggest that training exclusively on arXiv papers does not enhance mathematical reasoning abilities. In fact, for both models, arXiv-only pretraining yields little to no measurable gains—or even slight regressions—on a range of mathematical evaluation benchmarks examined in this work. These benchmarks span quantitative reasoning datasets such as GSM8K and MATH (refer to Table \ref{tab:arxiv-cot}), multiple-choice assessments including MMLU-STEM (Table \ref{tab:arxiv-cot}), and formal mathematics problems exemplified by miniF2F (Table \ref{tab:arxiv_atp}).

However, this conclusion has its limitations and should be taken with a grain of salt. We have not yet studied:

\begin{itemize}[topsep=0pt]
    \item The influence of arXiv tokens on other mathematics-oriented tasks excluded from this study, for example, transforming formal theorems or proofs into their informal counterparts;
    \item The consequences of incorporating arXiv tokens alongside different forms of data;
    \item The extent to which the advantages offered by arXiv papers might appear at increased model sizes.
\end{itemize}

Thus, further exploration is required, which we leave for future studies.

\subsection{Insights of Reinforcement Learning}
\subsubsection{Towards to a Unified Paradigm} Here, we present a comprehensive framework for examining a variety of training strategies including SFT, RFT, DPO, PPO, and GRPO. Additionally, we carry out experimental analyses to investigate the key components of this overarching framework. Typically, the gradient of a particular training method with respect to its parameter $\theta$ takes the form:

\begin{equation}
    \nabla_{\theta}\mathcal{J}_{\textcolor{red}{\mathcal{A}}}(\theta) = \mathbb{E}[\underbrace{(q,o) \sim \textcolor{red}{\mathcal{D}}}_{Data \ Source}]\left( \frac{1}{|o|} \sum_{t=1}^{|o|}  \underbrace{GC_{{\mathcal{A}}}(q, o, t, \textcolor{red}{\pi_{{rf}}})}_{Gradient \ Coefficient}  \nabla_{\theta}\log \pi_{\theta}(o_t | q, o_{<t})\right).
\label{eq:objective}
\end{equation}

There exist three key components: 1) \textit{Data Source $\mathcal{D}$}, which determines the training data; 2) \textit{Reward Function $\pi_{{rf}}$}, which is the source of the training reward signal; 3) \textit{Algorithm $\mathcal{A}$}: which processes the training data and the reward signal to the gradient coefficient $GC$ that determines the magnitude of the penalty or reinforcement for the data. We analyze several representative methods based on such a unified paradigm:

\begin{itemize}
    \item \textbf{Supervised Fine-tuning (SFT)}: In SFT, a pretrained model undergoes additional training using SFT data that has been curated by human annotators.
    \item \textbf{Rejection Sampling Fine-tuning (RFT)}: With RFT, the model trained with SFT is further updated by leveraging sampled responses to SFT prompts, with these outputs filtered and retained according to the accuracy of their answers.
    \item \textbf{Direct Preference Optimization (DPO)}: DPO advances the SFT model through continued fine-tuning using an expanded set of sampled outputs and employs a pairwise DPO loss to guide optimization.
    \item \textbf{Online Rejection Sampling Fine-tuning (Online RFT)}: Distinct from RFT, Online RFT sets the initial policy model from the SFT checkpoint but continually adapts it by incorporating newly generated outputs sampled on-the-fly from the evolving policy model itself.
    \item \textbf{PPO/GRPO}: For PPO/GRPO, the policy model is initialized from the SFT-trained model and enhanced by applying reinforcement learning using samples produced from the latest iteration of the policy model.
\end{itemize}

Table \ref{tab:data_policy} presents an overview of the elements constituting these approaches. For an in-depth explanation of the derivation procedure, consult Appendix \ref{app:analysis-rl}.

\paragraph{Observation about Data Source} The data source can be grouped into two main types: online sampling and offline sampling. In online sampling, the training data originates from real-time outputs generated by the current training policy model during its ongoing exploration. In contrast, offline sampling refers to data that has been collected from the sampling outputs of the initial SFT model prior to the commencement of further training. Methods such as RFT and DPO employ the offline sampling approach, whereas Online RFT and GRPO utilize online sampling.

As illustrated in Figure \ref{fig:rl-analysis}, our results indicate that Online RFT delivers markedly better performance than RFT across two benchmark tasks. Notably, Online RFT performs similarly to RFT at the beginning of training but achieves a clear performance lead as training progresses, thereby highlighting the benefits of the online training approach. This outcome is expected: during the initial training phase, the actor remains quite similar to the SFT model, so the data generated show only subtle variations. In contrast, during later stages, as the actor diverges from the SFT model, the sampled data reflect more pronounced distinctions, allowing real-time data collection to provide more substantial benefits.

\paragraph{Observation about Gradient Coefficient} The mechanism by which the algorithm updates model parameters involves translating the reward signal into the gradient coefficient. In our experimental setup, the reward function is bifurcated into `Rule' and `Model' categories. The `Rule' component evaluates response quality through assessment of answer accuracy, whereas the `Model' aspect involves training a reward model tasked with scoring each response, utilizing data labeled according to the rule-based judgments. As illustrated by Equations \ref{eq:GC-RFT} and \ref{eq:GC-GRPO}, there is a significant distinction between GRPO and Online RFT. Specifically, GRPO dynamically modulates its gradient coefficient in alignment with the reward values generated by the reward model, enabling nuanced differentiation in how varying levels of response quality are either reinforced or penalized. Conversely, Online RFT does not incorporate this level of granularity: it refrains from penalizing incorrect answers and applies the same reinforcement to all correct answers, lacking reward-dependent adjustment.

As illustrated in Figure \ref{fig:rl-analysis}, GRPO outperforms online RFT, emphasizing the advantage of modifying both positive and negative gradient coefficients. Moreover, GRPO+PS achieves better results than GRPO+OS, demonstrating the effectiveness of employing fine-grained, step-aware gradient coefficients. Additionally, we investigate the use of iterative RL by running two iterations in our experiments. Figure \ref{fig:iter-rl} reveals that iterative RL leads to substantial performance gains, particularly during the initial iteration.

\subsubsection{Why RL Works?} In this study, we employ reinforcement learning using a subset of instruction tuning data and observe substantial gains over the baseline instruction-tuned model. To shed light on the mechanisms behind reinforcement learning's effectiveness, we compare the Pass@K and Maj@K accuracies for both the Instruct and RL models across two benchmarks. Figure \ref{fig:maj-pass} illustrates that RL consistently improves Maj@K performance, while Pass@K remains largely unchanged. These results suggest that RL primarily strengthens the model’s performance by producing a more reliable output distribution, meaning \textbf{the observed benefit arises from elevating the likelihood of correct answers within the TopK outputs, instead of directly enhancing intrinsic model capability.} In a related vein, \citep{wang2023making} highlighted a \textbf{misalignment problem} for reasoning tasks in SFT models, and demonstrated that their reasoning abilities could be improved via various preference alignment methods \citep{yuan2023rrhf,song2023preference,wang2023making}.

\subsubsection{How to Achieve More Effective RL?}
\label{sec:effec_rl}
Our findings indicate that RL is highly effective for mathematical reasoning tasks. Additionally, we introduce a comprehensive framework to interpret various notable training strategies. In this framework, each method can be classified as either a direct application or an abstraction of RL methods. As highlighted in Equation \ref{eq:objective}, three essential elements emerge: Data Source, Algorithm, and Reward Function. We discuss several avenues for future exploration pertaining to these components.

\paragraph{Data Source} The foundation of all training approaches lies in the underlying data source. Within the realm of RL, this term denotes the collection of unlabeled questions paired with outputs that have been generated by the policy model. Our methodology restricts itself to employing questions taken from the instruction tuning phase, along with straightforward nucleus sampling as the output generation approach. We speculate that this limitation may contribute to the observed gains being confined to Maj@K metrics within our RL pipeline. Future investigations will focus on applying our RL methods to question prompts that fall outside the initial distribution, as well as integrating \textbf{advanced sampling (decoding) strategies}—such as those employing tree-search algorithms \citep{yao2023tree}. Furthermore, the application of \textbf{efficient inference techniques} \citep{xia-etal-2023-speculative,leviathan2023fast,kwon2023efficient,xia2024unlocking}, which greatly influence how effectively policy models can explore, represents another critical avenue for enhancing performance.

\paragraph{Algorithms} Algorithms utilize both the data and the reward signal to compute gradient coefficients, thereby guiding updates to the model parameters. As indicated in Equation \ref{eq:objective}, most current methods fundamentally \textbf{TRUST} the reward function's output to modulate the conditional probability associated with specific tokens, either increasing or decreasing it. Nevertheless, there is no guarantee that the reward signal is consistently accurate, particularly in highly intricate tasks. A case in point is the PRM800K dataset \citep{lightman2023let}, which, despite being meticulously labeled by expert annotators, still contains around 20\% annotation errors\footnote{\url{https://github.com/openai/prm800k/issues/12\#issuecomment-1728491852}}. In response, our focus will shift towards reinforcement learning algorithms designed to maintain robustness in the face of noisy reward signals. We posit that such \textbf{WEAK-TO-STRONG} alignment frameworks \citep{burns2023weak} have the potential to profoundly transform learning algorithm paradigms.

\paragraph{Reward Function} The reward function provides the core learning signal during training. In RL, this function typically corresponds to a neural-based reward model. We identify three critical avenues for advancing reward models: 1) \textbf{Improving the generalization capacity of the reward model.} The reward model needs to successfully generalize to handle questions outside its training distribution and accommodate more sophisticated decoding outcomes; without such ability, reinforcement learning risks simply maintaining the distributional properties of LLMs, rather than contributing to substantial capability enhancement; 2) \textbf{Capturing and expressing the uncertainty inherent to the reward model.} Characterizing this uncertainty can serve as an essential interface between imperfect reward models and learning algorithms that seek to bridge performance gaps; 3) \textbf{Constructing high-quality process reward models with efficiency} in order to deliver detailed, fine-grained feedback signals throughout the reasoning process \citep{lightman2023let,wang2023math}.

\section{Conclusion, Limitation, and Future Work}

This work introduces DeepSeekMath, a model that surpasses all existing open-source systems on the MATH benchmark at the competition level, and narrows the performance gap with proprietary models. DeepSeekMath is initialized from DeepSeek-Coder-v1.5 7B and is continuously trained on 500B tokens; notably, 120B math-focused tokens derived from Common Crawl form a major component of the dataset. Through a comprehensive ablation analysis, we reveal that web-scraped data holds considerable promise for curating high-quality mathematical corpora, whereas arXiv data may contribute less than initially anticipated. Additionally, we propose Group Relative Policy Optimization (GRPO), an adaptation of Proximal Policy Optimization (PPO), which effectively enhances mathematical reasoning while reducing memory requirements. Experimental findings indicate that GRPO continues to deliver performance gains even when DeepSeekMath-Instruct 7B achieves strong benchmark results. Finally, we offer a cohesive framework to interpret various methodologies and highlight several avenues for advancing reinforcement learning efficiency in future research.

While DeepSeekMath~demonstrates strong performance on benchmarks assessing quantitative reasoning, its abilities in the areas of geometry and theorem proving remain inferior to those of closed-source models. For example, during our preliminary evaluation, DeepSeekMath~was unable to address problems involving triangles and ellipses, suggesting the presence of data selection bias during pre-training and fine-tuning stages. Additionally, due to limitations in model size, DeepSeekMath~lags behind GPT-4 in few-shot learning scenarios. GPT-4 benefits from few-shot input by enhancing its results, whereas DeepSeekMath~shows comparable outcomes regardless of whether zero-shot or few-shot methodologies are applied. Moving forward, we plan to further refine our engineered data selection pipeline in order to create a more robust pre-training dataset. Furthermore, we aim to investigate new avenues (see Section \ref{sec:effec_rl}) to enable more efficient reinforcement learning for LLMs.

\end{document}